\chapter*{Kết luận}
\addcontentsline{toc}{chapter}{Kết luận}

Đề tài đã tập trung giải quyết bài toán xây dựng một ứng dụng di động hỗ trợ tập gym và theo dõi sức khỏe cá nhân hóa trên nền tảng Android. Trên cơ sở phân tích nhu cầu thực tế của người dùng phổ thông, đặc biệt là nhóm người mới bắt đầu hoặc chưa có khả năng tự xây dựng lộ trình tập luyện, hệ thống Fitty được thiết kế để không chỉ cung cấp một danh sách bài tập đơn lẻ mà còn tạo ra một luồng trải nghiệm hoàn chỉnh gồm xác thực truy cập, onboarding, sinh kế hoạch tập khởi đầu, tra cứu bài tập, hỗ trợ thực hiện buổi tập, theo dõi tiến độ cơ thể, ghi nhận dinh dưỡng và duy trì tương tác qua thông báo.

Về mặt kỹ thuật, hệ thống đã hiện thực được một kiến trúc nhiều tầng tương đối rõ ràng với giao diện Jetpack Compose, điều hướng bằng Navigation Compose, quản lý phụ thuộc bằng Hilt, lưu trữ cục bộ bằng Room và DataStore, cùng các dịch vụ phía sau dựa trên Firebase. Một trong những điểm nổi bật của hệ thống là việc áp dụng mô hình ngoại tuyến ưu tiên cho phân hệ bài tập, qua đó tách dữ liệu mô tả và media, đồng bộ dữ liệu từ xa theo cơ chế trung gian và sử dụng kho dữ liệu cục bộ như nguồn phục vụ ổn định cho giao diện. Bên cạnh đó, quy trình khởi tạo hồ sơ gắn với việc hình thành kế hoạch tập khởi đầu, quy trình hoàn thành buổi tập gắn với cập nhật tracking và việc lưu nhật ký bữa ăn, body scan là các điểm thể hiện rõ tính liên kết giữa các phân hệ của hệ thống.

Qua quá trình phân tích, thiết kế, cài đặt và kiểm thử, có thể đánh giá rằng hệ thống đã đáp ứng phần lớn những chức năng quan trọng của một ứng dụng hỗ trợ gymer theo hướng cá nhân hóa. Các phân hệ cốt lõi như xác thực, khởi tạo hồ sơ, kế hoạch tập, nội dung bài tập, buổi tập, theo dõi tiến độ, dinh dưỡng, hồ sơ người dùng, thông báo cục bộ và thông báo đẩy đều đã được hiện thực ở mức có thể vận hành như một nguyên mẫu chức năng có cấu trúc rõ ràng. Ngoài ra, việc có sẵn các kiểm thử đơn vị cho nhiều luồng nghiệp vụ quan trọng cũng là cơ sở tốt để tiếp tục mở rộng và nâng cao chất lượng hệ thống.

Tuy nhiên, hệ thống vẫn còn một số điểm cần tiếp tục hoàn thiện trong tương lai. Một số chỉ số theo dõi hiện vẫn được ước lượng từ dữ liệu sẵn có thay vì được hỗ trợ hoàn toàn từ một mô hình nghiệp vụ sâu hơn. Phần hướng dẫn ``động tác chuẩn'' hiện mới được hiện thực mạnh ở mức nội dung mô tả, media và gợi ý thực hiện, chưa phải là cơ chế phân tích chuyển động thời gian thực. Tương tự, phân tích dinh dưỡng và thể trạng mới đóng vai trò hỗ trợ ban đầu, chưa thay thế được đánh giá chuyên môn. Chất lượng trải nghiệm ở các màn hình theo dõi tiến độ, meal tracking và buổi tập nhanh còn có thể nâng cấp thêm nếu dữ liệu động và các cấu hình lưu trữ được làm giàu hơn. Bên cạnh đó, việc tăng cường kiểm thử tích hợp trên thiết bị thật và tiếp tục chuẩn hóa nội dung quản lý từ Firebase sẽ giúp hệ thống tiến gần hơn tới một sản phẩm có mức hoàn thiện cao.

Trong các hướng phát triển tiếp theo, hệ thống có thể được mở rộng theo nhiều hướng như cá nhân hóa kế hoạch tập sâu hơn, bổ sung biểu đồ theo dõi trực quan, nâng cấp phân tích dinh dưỡng và thể trạng, xây dựng meal planning theo mục tiêu tăng cơ hoặc giảm mỡ, tăng cường đồng bộ hóa nội dung hướng dẫn, tích hợp cảm biến hoặc dữ liệu sức khỏe từ thiết bị ngoài và phát triển các cơ chế đánh giá tư thế ở mức thông minh hơn. Với kiến trúc hiện tại, Fitty đã có nền tảng phù hợp để tiếp tục được phát triển thành một hệ thống hoàn chỉnh hơn trong phạm vi sản phẩm thực tế.

Tóm lại, đề tài đã đạt được mục tiêu xây dựng một hệ thống ứng dụng di động hỗ trợ tập luyện và theo dõi sức khỏe cá nhân hóa với kiến trúc hợp lý, phạm vi chức năng rõ ràng và khả năng mở rộng tốt. Các kết quả đạt được từ phân tích, thiết kế và cài đặt hệ thống là cơ sở có giá trị cho các bước phát triển tiếp theo của ứng dụng cũng như cho việc nghiên cứu và hiện thực các bài toán tương tự trong lĩnh vực ứng dụng sức khỏe trên thiết bị di động.
